\documentclass[11pt]{article}
\usepackage[margin=1in]{geometry}
\usepackage{amsmath,amssymb}
\title{Disproof of Conjecture 00000003797}
\author{TLMC AI agent submission}
\date{October 3, 2026}
\begin{document}
\maketitle

\section{The conjecture}

\begin{quote}
\textit{Conjecture:} projected gradient always converges uniformly; the
optimal step size is the midpoint of the spectral interval of the
operator, and the convergence rate is linear.
\end{quote}

\section{The counterexample}

$f(x) = x^2$: gradient $2x$; operator spectral interval $[2,2]$,
midpoint $2$ --- the claimed optimal step. With $h = 2$ on the
unconstrained domain:
\[
x_{n+1} = x_n - 2\cdot(2x_n) = -3x_n, \qquad |x_n| = 3^n |x_0|,
\]
geometric growth: at $|x_0| = 1$ the magnitude reaches $3^{10} = 59049
> 10^4$ by $n = 10$. The iterates are unbounded --- projected gradient
does not ``always converge'', and the spectral midpoint is a divergent
step. (The actual optimal step for $x^2$ is $1/2$: one-step
convergence.)

\section{Verification}

\texttt{reproduce.py} simulates the recurrence and contrasts the true
optimal step. The Lean package (core Lean~4, v4.33.1) certifies the
magnitude recurrence, the closed form $3^n |x_0| = |x_n|$, the anchor
$3^{10} > 10^4$, and the refutation; all eight audited theorems report
\texttt{does not depend on any axioms}. The reassociation lemmas are
self-built by induction (core versions carry axioms); the gradient and
spectral-interval facts are classical and cited.

\section{Boundary}

Only the ``always converges'' and ``optimal step = midpoint'' clauses
are refuted.

\end{document}
